\textbf{4-1}\quad 质量为 $\mu$ 的粒子在三维球方势阱
$$
V (r) = \left\{ \begin{array}{l l} 0, & r > a \\ - V _ {0}, & r <   a \end{array} \right. \quad (V _ {0} > 0)
$$

中运动. 求存在 s 波束缚态的条件.

\vfill

\newpage
\vfill

\noindent
\textbf{4-2}\quad 粒子处于三维球壳势阱 $V(r) = -g\delta(r - a)$ 中, 求存在束缚态的条件.

\vfill

\newpage
\vfill

\noindent
\textbf{4-3}\quad 质量为 $\mu$ 的粒子在势场 $V(r) = \left\{ \begin{array}{ll}0, & a < r < b\\ \infty , & r\leqslant a,r\geqslant b \end{array} \right.$ 中运动，求 $l = 0$ 的定态能量和定态波函数.

\vfill

\newpage
\vfill

\noindent
\textbf{4-4}\quad 设粒子的定态波函数 $\psi(r) = A\mathrm{e}^{-r / a}$ , 其中 $A$ 与 $a$ 是常数. 已知 $r \to \infty$ , $V(r) = 0$ . 求定态能量 $E$ 和势能 $V(r)$ .

\vfill

\newpage
\vfill

\noindent
\textbf{4-5}\quad 一个质量为 $\mu$ 带有电荷 $q$ 的粒子被限制在 $xy$ 平面内的半径为 $a$ 的圆周上运动. 在 $z$ 轴方向加上强度为 $B$ 的均匀磁场, 求粒子的基态能量和基态波函数的表达式, 证明基态能量是 $B$ 的周期函数, 并给出周期来. 已知在柱坐标系中
\begin{equation*}
\nabla^ {2} = \frac {1}{\rho} \frac {\partial}{\partial \rho} \rho \frac {\partial}{\partial \rho} + \frac {1}{\rho^ {2}} \frac {\partial^ {2}}{\partial \varphi^ {2}} + \frac {\partial^ {2}}{\partial z ^ {2}}
\end{equation*}

\vfill

\newpage
\vfill

\noindent
\textbf{4-6}\quad 质量为 $\mu$ 电荷为 q 的粒子在均匀磁场 B = Bk 中运动, 求定态能量与波函数.

\vfill

\newpage
\vfill

\noindent
\textbf{4-7}\quad 质量为 $\mu$ 电荷为 q 的粒子在方向互相垂直的均匀电场 E 和均匀磁场 B 中运动, 求定态能量与波函数.

\vfill

\newpage
\vfill

\noindent
\textbf{4-8}\quad (1)已知带有电荷 $q$ 的粒子在磁场 $\pmb{B}$ 与势场 $V(r)$ 中运动, 求带电粒子速度分量算符的对易关系 $[\hat{\nu}_x, \hat{\nu}_y]$ , $[\hat{\nu}_y, \hat{\nu}_z]$ , $[\hat{\nu}_z, \hat{\nu}_x]$ 的表达式. 

(2)质量为 $\mu$ 带有电荷 $q$ 的粒子在磁场中的哈密顿量为 $\hat{H} = \frac{1}{2\mu}\left(\hat{p} - \frac{q}{c} A\right)^2$ . 请问在什么条件下, 它可以写成如下形式: $\hat{H}' = \frac{1}{2\mu}\hat{p}^2 - \frac{q}{\mu c} A \cdot \hat{p} + \frac{q^2}{2\mu c^2} A^2$ . 

(3)设 $A = A_0 \cos \omega t$ ( $A_0$ 为常数矢量), 略去 $\hat{H}'$ 中的 $A^2$ 项, 求出与 $\hat{H}'$ 相应的薛定谔方程的解.

\vfill

\newpage
\vfill

\noindent
\textbf{4-9}\quad 处于基态的类氢原子经 $\beta$ 衰变, 核电荷数突然由 $Z$ 变为 $Z + 1$ , 求原子处于 $2s$ 态的概率. 已知类氢原子定态波函数为
$$
\psi_ {1 0 0} (Z, \boldsymbol {r}) = \frac {1}{\sqrt {\pi}} \left(\frac {Z}{a}\right) ^ {3 / 2} \mathrm{e} ^ {- Z r / a}, \psi_ {2 0 0} (Z, \boldsymbol {r}) = \frac {1}{\sqrt {\pi}} \left(\frac {Z}{2 a}\right) ^ {3 / 2} \left(1 - \frac {Z r}{2 a}\right) \mathrm{e} ^ {- Z r / 2 a}
$$

\vfill

\newpage
\vfill

\noindent
\textbf{4-10}\quad 氢原子处于基态. 假定库仑作用在 $t = 0$ 时突然消失, 电子离开原子像自由电子那样运动. 试求 $t > 0$ 时测量电子动量大小在 $p \sim p + \mathrm{d}p$ 内的概率.

\vfill

\newpage
\vfill

\noindent
\textbf{4-11}\quad 在半径为 $R$ 的硬刚球内, 有一质量为 $\mu$ 的粒子. 

(1) 求粒子的基态能量和波函数; 

(2) 如 $t < 0$ 时粒子处于基态, $t = 0$ 时将这硬刚球的半径扩展到原来的 2 倍, 求扩展后粒子仍处于基态的概率.

\vfill

\newpage
\vfill

\noindent
\textbf{4-12}\quad 一粒子被束缚在半径为 $R$ 的刚球盒内. 求处于基态的粒子对盒壁的压力与压强.

\vfill

\newpage
\vfill

\noindent
\textbf{4-13}\quad 在势场 $V(\pmb{r})$ 中粒子处于定态, 证明粒子动能 $\hat{T} = \hat{p}^2 / 2\mu$ 的平均值为 $\overline{T} = \frac{1}{2} (\overline{\pmb{r} \cdot \nabla V})$ . 如果 $V(\pmb{r})$ 是 $\pmb{r}$ 的 $\nu$ 次齐次函数, 证明 $\overline{T} = \frac{\nu}{2}\overline{V}$ , 并利用此式, 计算氢原子基态的 $\overline{T}$ .

\vfill

\newpage
\vfill

\noindent
\textbf{4-14}\quad 势能 $V = -\frac{Ze^2}{r}$ 的类氢原子处于 $\psi_{nlm}$ 态. 试计算 $\frac{1}{r}$ 的平均值 $\langle nlm|\frac{1}{r}|nlm\rangle$ .

\vfill

\newpage
\vfill

\noindent
\textbf{4-15}\quad 势能 $V = -\frac{Ze^2}{r}$ 的类氢原子处于 $\psi_{nlm}$ 态. 试计算 $\frac{1}{r^2}$ 的平均值 $\left\langle nlm \mid \frac{1}{r^2} \mid nlm \right\rangle$

\vfill

\newpage
\vfill

\noindent
\textbf{4-16}\quad 设一粒子在中心力场 $V(r)$ 中运动, 定义径向动量
\begin{equation}\tag{1}
\hat {p} _ {r} = \frac {1}{2} \left(\frac {\boldsymbol {r}}{r} \cdot \hat {\boldsymbol {p}} + \hat {\boldsymbol {p}} \cdot \frac {\boldsymbol {r}}{r}\right)
\end{equation}

(1) 证明
\begin{equation}\tag{2}
\hat {p} _ {r} = - \mathrm{i} \hbar \left(\frac {\partial}{\partial r} + \frac {1}{r}\right)
\end{equation}
\begin{equation}\tag{3}
\hat {H} = \frac {\hat {p} _ {r} ^ {2}}{2 \mu} + \frac {\tilde {L} ^ {2}}{2 \mu r ^ {2}} + V (r)
\end{equation}

(2) 计算对易关系式 $\left[\frac{\partial}{\partial r}, \hat{H}\right]$ ; 

(3) 当粒子处于某一束缚定态 $\psi_{nlm}(\boldsymbol{r}) = R_{nl}(r)Y_{lm}(\theta, \varphi)$ 时, 证明
\begin{equation}\tag{4}
\left\langle \frac {\partial V}{\partial r} \right\rangle_ {n l m} - \frac {l (l + 1) \hbar^ {2}}{\mu} \left\langle \frac {1}{r ^ {3}} \right\rangle_ {n l m} = \frac {\hbar^ {2}}{2 \mu} \left| R _ {n l} (0) \right| ^ {2}
\end{equation}

(4) 设 $V = -\frac{Ze^2}{r}$ , 证明在束缚定态 $\psi_{nlm}(\pmb{r})$ 上,
\begin{equation}\tag{5}
\left\langle \frac {1}{r ^ {3}} \right\rangle_ {n l m} = \frac {Z}{l (l + 1) a} \left\langle \frac {1}{r ^ {2}} \right\rangle_ {n l m}
\end{equation}

\vfill

\newpage
\vfill

\noindent
\textbf{4-17}\quad 氢原子处于基态 $\psi(r)=\frac{1}{\sqrt{\pi a^{3}}}\mathrm{e}^{-r/a}$ ，求 r 的平均值及动量的概率分布函数 $W(p)$ .

\vfill

\newpage
\vfill

\noindent
\textbf{4-18}\quad 氢原子处于基态 $\psi(r)=\frac{1}{\sqrt{\pi a^{3}}}\mathrm{e}^{-r/a}$ ，计算 $\Delta x\Delta p_{x}$ ，检验测不准关系.

\vfill

\newpage
\vfill

\noindent
\textbf{4-19}\quad 固有长度为 $r_0$ 的平面转子处于状态 $\psi_m(\varphi) = \sqrt{\frac{1}{2\pi}}\mathrm{e}^{\mathrm{i}m\varphi}$ ( $m = 0, \pm 1, \pm 2, \dots$ ), 其中平面极角 $\varphi$ 与坐标 $x$ 和 $y$ 的关系是 $x = r_0\cos \varphi, y = r_0\sin \varphi$ . 求坐标 $x$ 和动量 $p_x$ 的不确定关系 $\Delta x\Delta p_x = ?$

\vfill

\newpage
\vfill

\noindent
\textbf{4-20}\quad 利用测不准关系估算氢原子基态能量.

\vfill

\newpage
\vfill

\noindent
\textbf{4-21}\quad 利用测不准关系, 估算质量为 $\mu$ 的粒子在如下势场中的基态能量: (1) $V(r) = kr(k > 0)$ ; (2) $V(r) = -\frac{\lambda}{r^{3/2}} (\lambda > 0)$ .

\vfill

\newpage
\vfill

\noindent
\textbf{4-22}\quad 原子核的线度为 $10^{-13} \mathrm{~cm}$ . 试用不确定原理估算核内质子的动能(以电子伏特为单位).

\vfill

\newpage
\vfill

\noindent
\textbf{4-23}\quad 一个质量为 $\mu$ 的粒子在对数势场 $V(r) = c\ln \frac{r}{r_0}$ 中运动, 其中 $c$ 与 $r_0$ 是同质量 $\mu$ 无关的常数. 

(1)证明在所有定态上均方速度相同, 求出这个均方速度; 

(2)证明任何两个定态能量之差同粒子的质量无关.

\vfill

\newpage
\vfill

\noindent
\textbf{4-24}\quad 两个质量均为 $m$ 的粒子, 通过三维球对称势 $V(r) = c \ln (r / r_0) (c > 0, r_0 > 0)$ 而束缚在一起, $r$ 为两粒子之间的距离. 已知它的第一激发态与基态的能量之差为 $\Delta E$ . 今有一个质量为 $m$ 的粒子与另一个质量为 $1840m$ 的粒子通过同一位势形成束缚态, 求这一系统第一激发态与基态的能量之差. 请说明理由, 并给以证明.

\vfill

\newpage
\vfill

\noindent
\textbf{4-25}\quad 设一微观粒子在中心力场 $V(r)$ 中运动, 且处于能量和轨道角动量的某一共同本征态. 

(1) 在球坐标系中写出能量本征函数的基本形式, 写出势能 $V(r)$ 在此态上的平均值 $\langle V\rangle$ 的表达式, 并最后表示成径向积分的形式; 

(2) 设 $V(r)$ 是 $r$ 的单调上升函数 (对任意 $r$ , $\mathrm{d}V / \mathrm{d}r > 0$ ), 证明对任意给定的 $r_0$ , 均有 $\int_0^{r_0}[V(r) - \langle V\rangle]R^2 (r)r^2\mathrm{d}r < 0$ , 其中 $R(r)$ 是径向波函数.

\vfill

\newpage
\vfill

\noindent
\textbf{4-26}\quad 在 $t = 0$ 时, 氢原子的波函数为
$$
\psi (r, 0) = \frac {1}{\sqrt {1 0}} \left(2 \psi_ {1 0 0} + \psi_ {2 1 0} + \sqrt {2} \psi_ {2 1 1} + \sqrt {3} \psi_ {2 1 - 1}\right)
$$
其中下标分别是量子数 $n, l, m$ 的值, 不考虑自旋. 

(1) 求体系的平均能量; 

(2) 在任意 $t$ 时刻体系处于 $l = 1, m = 1$ 的态的概率是多少? 

(3) 在任意 $t$ 时刻体系处于 $m = 0$ 的态的概率是多少? 

(4) 写出任意 $t$ 时刻体系的波函数 $\psi(r, t)$ .

\vfill

\newpage
\vfill

\noindent
\textbf{4-27}\quad 质量为 $\mu$ 电荷为 $q$ 的粒子在均匀恒定磁场中运动, 取不对称规范: $A_x = -By$ , $A_y = A_z = 0$ , $B$ 为磁场大小, 则可知 $y_0 = -c\hat{p}_x / qB$ 是守恒量. 证明下面的量: $x_0 = x + (c\hat{p}_y / qB)$ 也是守恒量, 它与 $y_0$ 是否可以同时被观测?

\vfill

\newpage
\vfill

\noindent
\textbf{4-28}\quad 质量为 $\mu$ 的粒子在中心力场 $V(r) = -\frac{\alpha}{r^s} (\alpha > 0)$ 中运动. 证明存在束缚态的条件是 $0 < s < 2$ .

\vfill

\newpage
\vfill

\noindent
\textbf{4-29}\quad 氘核是由质子和中子组成的唯一的束缚态. 实验测定氘核的结合能为 2.23MeV. 设质子和中子的作用力势可近似表示为 $V(r) = \begin{cases} -V_0, & r < a \\ 0, & r > a \end{cases}$ , 求作用力程 $a$ 和作用强度 $V_0$ 之间的关系式.

\vfill

\newpage
\vfill

\noindent
\textbf{4-30}\quad 双原子分子中两原子之间的作用力势可表示为
$$
V (r) = V _ {0} \left[ 1 - \mathrm{e} ^ {- (r - R) / a} \right] ^ {2} - V _ {0}
$$

其中 $r$ 为两原子之间的距离， $R$ 与 $a$ 为正的常数， $R > a$ ，且 $\mathbf{e}^{R / a} \gg 1$ 。当 $r \to 0$ 时， $V(r) \to \infty$ ；随 $r$ 增大， $V(r)$ 迅速变小；当 $r = R$ 时， $V(r)$ 取最小值 $-V_0$ ；当 $r \to \infty$ 时， $V(r) \to 0$ 。求轨道角动量 $l = 0$ 的束缚定态能量。

\vfill

\newpage
\vfill

\noindent
\textbf{4-31}\quad 粒子在中心力场 $V(r) = \frac{A}{r^2} -\frac{B}{r}$ 中运动, 其中 $A$ 与 $B$ 为正实数. 求定态能量和波函数. (提示: 将定态方程化为氢原子定态方程)

\vfill

\newpage
\vfill

\noindent
\textbf{4-32}\quad 粒子在势场 $V(r) = -V_0\mathrm{e}^{-r / a}$ 中运动, 其中 $V_{0}$ 与 $a$ 为正实数. 求存在束缚态的条件.

\vfill

\newpage
\vfill

\noindent
\textbf{4-33}\quad 设氢原子处于 $\psi(\boldsymbol{r}) = R_{21}(r) Y_{1-1}(\theta, \varphi)$ 态, 其中

$$
R _ {2 1} (r) = \frac {1}{2 \sqrt {6} a ^ {3 / 2}} \frac {r}{a} \mathrm{e} ^ {- r / 2 a}, Y _ {1 - 1} (\theta , \varphi) = \sqrt {\frac {3}{8 \pi}} \sin \theta \mathrm{e} ^ {- \mathrm{i} \varphi}
$$

(1) 求势能 $V = -e^{2} / r$ 的平均值 $\left\langle V\right\rangle$ ；

(2) 求轨道角动量 $\hat{L}_z \hat{L}_x^2$ 的平均值 $\left\langle \hat{L}_z \hat{L}_x^2 \right\rangle$

\vfill

\newpage
\vfill

\noindent
\textbf{4-34}\quad 质量为 $\mu$ 电荷为 $q$ 的粒子在三维各向同性谐振子场 $V(r) = \frac{1}{2}\mu \omega^2 r^2$ 中运动, 同时受到一个沿 $x$ 方向的均匀电场 $E = E_0i$ 的作用. 求粒子的能量本征值和第一激发态的简并度. 此时轨道角动量是否守恒? 如回答是, 写出此时守恒量的表达式.

\vfill

\newpage
\vfill

\noindent
\textbf{4-35}\quad 二维平面有一个以原点为圆心, 半径为 $R$ 的光滑刚性圆环, 上面套着两颗质量为 $m$ , 电荷为 $q$ 的"量子珍珠". "量子珍珠"可以无摩擦地绕环滑动, 相互之间存在库仑位势 $V = q^2 / r, r$ 为两颗珍珠之间的距离, 两颗珍珠各自的方位角(或极角)为 $\alpha, \beta$ . 

(1) 利用 $\alpha, \beta$ 写出这一珍珠体系的定态薛定谔方程; 

(2) 令 $\lambda = (\alpha + \beta) / 2, \theta = \alpha - \beta$ , 证明定态方程可以利用 $\lambda, \theta$ 分离变量, 从而这一体系可以分解为关于 $\lambda$ 的"质心"运动和关于 $\theta$ 的相对运动; 

(3) 求解关于 $\lambda$ 的"质心"运动; 

(4) 当 $q^2$ 非常大时, 求解或估算关于 $\theta$ 的相对运动的基态与最低激发态的本征能量.

\vfill

\newpage
\vfill

\noindent
\textbf{4-36}\quad 粒子被限制在无限长圆筒内运动, 圆筒半径为 $a$, 求粒子的能量.

\vfill

\newpage